\section{Architecture matérielle et chaînes RF d’un radar SAR embarqué}\label{SEC:Architecture matérielle et chaînes RF d’un radar SAR embarqué}

\subsection{Architecture fonctionnelle globale}

Chaîne émission → antenne → propagation → réception → base bande.

Différences entre chaîne RF SAR pulsé et FMCW-SAR.


\subsection{Sous-systèmes RF}

Source RF / Synthétiseur : VCO, PLL, DDS, stabilité de phase critique pour SAR.

Émetteur : modulateur FMCW ou impulsionnel, amplificateur de puissance (PA).

Récepteur : LNA, mélangeurs, filtres, amplificateurs IF, démodulateur IQ.

Conversion fréquence / Base bande : downconversion, ADC haute dynamique.


\subsection{Antenne et front-end RF}

Antennes à ouverture synthétique : antennes linéaires, réseaux patch, antennes conformes.

Contraintes d’intégration (taille, masse, diagramme de rayonnement, polarisation).

Antennes actives vs passives, antennes MIMO pour SAR.


\subsection{Gestion des données}

Débits très élevés (plusieurs centaines de Mo/s).

Stockage embarqué (SSD industriels) ou transfert vers station au sol.


\subsection{Alimentation et contraintes embarquées}

Gestion énergétique, rendement RF.

Dissipation thermique, miniaturisation et fiabilité.